\section{Build Optimization}
\label{sec:build}

Datastore's binary is built with three composable techniques: profile-guided
optimization (PGO), ThinLTO, and the BOLT post-link optimizer. We describe the
pipeline and the constraints that order it; Section~\ref{sec:eval} reports the
effect of each stage.

\subsection{Pipeline}

The canonical optimized build (\texttt{utils/pgo-build/build-optimized.sh})
proceeds in three phases:

\begin{enumerate}[leftmargin=1.5em]
  \item \textbf{Instrumented build.} Compile with clang IR instrumentation
  (\texttt{-fprofile-generate}). Because IR-level PGO instrumentation is
  incompatible with ThinLTO, this phase runs with LTO disabled.
  \item \textbf{Profile collection.} Run a representative OLAP workload ---
  inserts into a \texttt{MergeTree} table followed by aggregation, top-$k$,
  filter, time-bucketing, join, window, array, and sub-query reads --- against
  the instrumented binary, then merge the raw counters with
  \texttt{llvm-profdata}.
  \item \textbf{Optimized build.} Recompile with \texttt{-fprofile-use} and
  ThinLTO enabled. Functions whose control-flow hash no longer matches the
  recorded profile (expected, since instrumentation ran without LTO) are
  demoted from \texttt{-Werror} to a warning so the build does not abort on
  benign profile staleness.
\end{enumerate}

The resulting binary is optionally passed to BOLT, which re-lays-out hot code
using a branch profile, separating hot and cold basic blocks and reordering
functions for instruction-cache locality. BOLT requires the linker to emit
relocations (\texttt{-Wl,--emit-relocs}), which we gate behind a build flag.

\subsection{A subtle thread-pool deadlock in instrumented collection}

A practical obstacle: running the \emph{instrumented} binary against a
large insert deadlocked. The \texttt{-fprofile-generate} runtime serializes
access to per-function counters, and this serialization interacts badly with
the engine's background merge thread pools under a high-volume
\texttt{MergeTree} insert. We resolve this by collecting the profile with
bounded threads (\texttt{max\_insert\_threads=1},
\texttt{background\_pool\_size=1}) over a smaller but representative row count.
PGO requires that the profile cover the hot \emph{code paths}, not a
production-scale data volume, so this does not degrade profile quality; it
completes in seconds rather than deadlocking.

\subsection{Profile representativeness and matched comparison}

A PGO build is only as good as its profile. For a fair comparison we hold the
profile fixed across the fork and upstream, and --- crucially --- we regenerate
the profile against the \emph{same} build configuration the binary will ship
with. An early result used a profile collected from an XRay-instrumented build
to optimize a non-XRay binary; the control-flow-hash mismatch on hot functions
silently discarded counts for exactly the aggregation kernels we cared about,
understating the aggregation result. The shipped recipe collects the profile
from the final, XRay-free configuration, eliminating the mismatch.
